\documentclass[british]{article}

\usepackage[margin=1.5in]{geometry}

\title{Mixed manifolds and mixed spaces}
\author{Adrien Douady}
\date{1960}

% Standard
\usepackage{amssymb,amsmath,amscd}
\usepackage{enumerate}
\usepackage{tikz-cd}
\usepackage{mathtools}
\usepackage{booktabs}


%Bibliography
\usepackage[backend=biber,maxnames=99,maxalphanames=4]{biblatex}
\addbibresource{SHC-13_1-2.bib}


% Colours
\usepackage{xcolor}


% Fonts
\usepackage{mathrsfs}
\usepackage{fouriernc}


% Things for Quarto/pandoc
\definecolor{quarto-callout-color-frame}{HTML}{000000}
\usepackage{calc,array}
\usepackage{longtable}
\providecommand{\tightlist}{%
\setlength{\itemsep}{0pt}\setlength{\parskip}{0pt}}
\newcounter{none}
\usepackage{graphicx}
\makeatletter
\newsavebox\pandoc@box
\newcommand*\pandocbounded[1]{% scales image to fit in text height/width
  \sbox\pandoc@box{#1}%
  \Gscale@div\@tempa{\textheight}{\dimexpr\ht\pandoc@box+\dp\pandoc@box\relax}%
  \Gscale@div\@tempb{\linewidth}{\wd\pandoc@box}%
  \ifdim\@tempb\p@<\@tempa\p@\let\@tempa\@tempb\fi% select the smaller of both
  \ifdim\@tempa\p@<\p@\scalebox{\@tempa}{\usebox\pandoc@box}%
  \else\usebox{\pandoc@box}%
  \fi%
}
% Set default figure placement to htbp
\def\fps@figure{htbp}
\makeatother
\usepackage{kvsetkeys}


% Header and footer
\usepackage{fancyhdr}
\usepackage{lastpage}
\usepackage{datetime2}
\pagestyle{fancy}
\fancypagestyle{plain}{}
\fancyhf{}
\renewcommand{\headrulewidth}{0pt}
\cfoot{\small\thepage\ of \pageref*{LastPage}}
\lfoot{\footnotesize Build date: \today}


% Theorem environments
\usepackage{amsthm}
\newenvironment{itenv}[1]
  {\phantomsection\par\smallskip\noindent\textbf{#1.}\itshape}
  {\par\smallskip}
\newenvironment{rmenv}[1]
  {\phantomsection\par\smallskip\noindent\textbf{#1.}\rmfamily}
  {\par\smallskip}


% Shortcuts
\newcommand{\oldpage}[1]{\marginpar{\footnotesize$\Big\vert$ \textit{p.~#1}}}


%% Document %%

% hyperref last, as is "good practice"

\usepackage{hyperref}
\hypersetup{colorlinks,linkcolor={purple!50!black},citecolor={purple!50!black},urlcolor={purple!80!black}}

\usepackage{embedall}
\begin{document}

\maketitle
\thispagestyle{fancy}

\setcounter{footnote}{0}


%% Content %%

\emph{This document is a translation into English of the following:}

\begin{quote}
Douady, A. ``Variétés et espaces mixtes''. \emph{Séminaire Henri Cartan}
\textbf{13 (1)} (1960--61), Talk no. 2.
\href{https://www.numdam.org/item/SHC_1960-1961__13_1_A1_0}{numdam.org/item/SHC\_1960-1961\_\_13\_1\_A1\_0}
\end{quote}

\emph{The translator (\href{https://thosgood.net}{Tim Hosgood}) takes
full responsibility for any errors introduced in this document, and
claims no rights to any of the mathematical content.}

\tableofcontents

\providecommand{\scr}[1]{{\mathscr{#1}}}
\renewcommand{\cal}[1]{{\mathcal{#1}}}
\renewcommand{\frak}[1]{{\mathfrak{#1}}}
\renewcommand{\geq}{\geqslant}
\renewcommand{\leq}{\leqslant}

\providecommand{\ZZ}{\mathbb{Z}}
\providecommand{\RR}{\mathbb{R}}
\providecommand{\CC}{\mathbb{C}}
\providecommand{\DD}{\mathrm{D}}
\providecommand{\HH}{\mathrm{H}}
\providecommand{\PP}{\mathbb{P}}
\providecommand{\dd}{\mathrm{d}}

\providecommand{\Hom}{\operatorname{Hom}}
\providecommand{\GL}{\operatorname{GL}}
\providecommand{\SL}{\operatorname{SL}}
\providecommand{\Ker}{\operatorname{Ker}}

\section*{I. Category of models}\label{two-section-I}
\addcontentsline{toc}{section}{I. Category of models}

\oldpage{2-01}

Let \(B\) be a topological space. We define the category
\(\mathscr{S}_B^n\) in the following manner: the objects of
\(\mathscr{S}_B^n\) are the open subsets of \(B\times\mathbb{C}^n\), and
a morphism \(f\colon U\to U'\) from an open subset
\(U\subset B\times\mathbb{C}^n\) to an open subset
\(U'\subset B\times\mathbb{C}^n\) is a continuous map
\(f\colon U\to U'\) satisfying the following two conditions:

\begin{enumerate}
\def\labelenumi{\arabic{enumi}.}
\tightlist
\item
  the diagram \[
     \begin{CD}
       U @>f>> U'
     \\@V{\pi_1}VV @VV{\pi_1}V
     \\B @= B
     \end{CD}
   \] commutes, where \(\pi_1\) denotes the projection of
  \(B\times\mathbb{C}^n\) to \(B\); and
\item
  for all \(x\in B\), the map \(f_x\colon U_x\to U'_x\) is holomorphic,
  where \[
     U_x = \{z\in\mathbb{C}^n \mid (x,z)\in U\}
   \] (and similarly for \(U'\)).
\end{enumerate}

If \(B\) is endowed with the structure of a \(\mathscr{C}^\infty\)
manifold (resp. an \(\mathbb{R}\)-analytic manifold, resp.
\(\mathbb{C}\)-analytic manifold), then we obtain a category
\(\mathscr{C}^\infty\mathscr{S}_B\) (resp. \(\mathbb{R}\mathscr{S}_B\),
resp. \(\mathbb{C}\mathscr{S}_B\)) by requiring the morphisms to be
\(\mathscr{C}^\infty\) (resp. \(\mathbb{R}\)-analytic, resp.
\(\mathbb{C}\)-analytic).

More generally, if \(f_1\colon B\to B'\) is a continuous map from one
topological space to another, then a \emph{morphism of
\(\mathscr{S}_{f_1}\)} is a continuous map \(f\) from an object \(U\) of
\(\mathscr{S}_B\) to an object \(U'\) of \(\mathscr{S}_{B'}\) such that

\begin{enumerate}
\def\labelenumi{\arabic{enumi}.}
\tightlist
\item
  the diagram \[
     \begin{CD}
       U @>f>> U'
     \\@V{\pi_1}VV @VV{\pi_1}V
     \\B @>>{f_1}> B'
     \end{CD}
   \] commutes; and
\item
  \(f_x\colon U_x\to U'_{f_1(x)}\) is holomorphic for all \(x\in B\).
\end{enumerate}

\oldpage{2-02}

If \(f_1\) is a \(\mathscr{C}^\infty\) map from one
\(\mathscr{C}^\infty\) manifold to another, then \(f\) will be a
morphism of \(\mathscr{C}^\infty\mathscr{S}_{f_1}\) if, further, it is a
\(\mathscr{C}^\infty\) map (resp. \ldots). We thus obtain, for every
category of topological spaces, a fibred category \(\mathscr{S}^n\)
(resp. \(\mathscr{C}^\infty\mathscr{S}^n\), resp. \ldots).

\section*{II. The definition of mixed spaces and mixed
varieties}\label{two-section-II}
\addcontentsline{toc}{section}{II. The definition of mixed spaces and
mixed varieties}

\subsection*{1. First definition}\label{two-section-II.1}
\addcontentsline{toc}{subsection}{1. First definition}

Let \(B\) and \(V\) be separated spaces, and let \(\pi\colon V\to B\) be
a continuous map. The structure of a \emph{mixed space} over \(B\) is
defined on \(V\) by a system of charts \(\varphi_i\colon U_i\to V\),
where the \((U_i)\) are objects of \(\mathscr{S}_B^n\); for each \(i\),
\(\varphi_i\) is a homeomorphism from \(U_i\) to an open subset of \(V\)
such that the diagram \[
  \begin{CD}
    U_i @>{\varphi_i}>> V
  \\@V{\pi_1}VV @VV{\pi}V
  \\B @= B
  \end{CD}
\] commutes; finally, for all \(i\) and all \(j\), the ``change of
chart'' \(\varphi_j^{-1}\circ\varphi_i\) is an isomorphism of
\(\mathscr{S}_B\) from an open subset of \(U_i\) to an open subset of
\(U_j\).

The structure thus defined is that of a
\emph{\((\mathscr{C}^0,\mathbb{C})\)-mixed space}. If \(B\) is a
\(\mathbb{C}\)-analytic space, and if the change of chart maps are all
\(\mathbb{C}\)-analytic, then we have a \emph{\(\mathbb{C}\)-analytic
mixed space}. In this case, \(V\) itself is a \(\mathbb{C}\)-analytic
space, and the fibres \(V_x=\pi^{-1}(x)\) are \(\mathbb{C}\)-analytic
sub-manifolds.

If \(B\) is a \(\mathscr{C}^\infty\) manifold (resp.
\(\mathbb{R}\)-analytic, resp. \(\mathbb{C}\)-analytic), and if the
change of chart maps are all \(\mathscr{C}^\infty\) (resp. \ldots), then
we have a \emph{\((\mathscr{C}^\infty,\mathbb{C})\)-mixed manifold}
(resp. \((\mathbb{R},\mathbb{C})\), resp. \((\mathbb{C},\mathbb{C})\)).
In this case, \(V\) itself is a manifold. Note that the notion of a
\((\mathbb{C},\mathbb{C})\)-mixed manifold, or a \(\mathbb{C}\)-analytic
mixed manifold, reduces to simply having a \(\mathbb{C}\)-analytic
manifold \(V\) endowed with a projection \(\pi\colon V\to B\) onto
another \(\mathbb{C}\)-analytic manifold such that \(\pi\) is of maximal
rank at every point.\footnote{\emph{{[}Trans.{]} The more common modern
  nomenclature is to simply call such an object a family of complex
  manifolds.}}

Let \(\pi\colon V\to B\) and \(\pi'\colon V'\to B'\) be mixed spaces,
and let \(f_1\colon B\to B'\) be a continuous (resp. \ldots) map. Then a
\emph{morphism from \(V\) to \(V'\) over \(f_1\)} is a continuous map
\(f\colon V\to V'\) such that the diagram \[
  \begin{CD}
    V @>f>> V'
  \\@V{\pi}VV @VV{\pi'}V
  \\B @>>{f_1}> B'
  \end{CD}
\] commutes, and such that, for any charts \(\varphi_i\colon U_i\to V\)
and \(\varphi'_j\colon U'_j\to V'\), the map
\({\varphi'_j}^{-1}\circ f\circ\varphi_i\) is a morphism of
\(\mathscr{S}_{f_1}\) (resp. \ldots) from an open subset of \(U_i\) to
\(U_j\). \oldpage{2-03}

\subsection*{2. An equivalent definition}\label{two-section-II.2}
\addcontentsline{toc}{subsection}{2. An equivalent definition}

We now give another way of defining mixed spaces, equivalent to the
above.

Given separated spaces \(B\) and \(V\), along with a continuous map
\(\pi\colon V\to B\), the structure of a \emph{pre-mixed space} consists
of the structure of a \(\mathbb{C}\)-analytic manifold on each fibre
\(V_x=\pi^{-1}(x)\). Given pre-mixed spaces \(\pi\colon V\to B\) and
\(\pi'\colon V'\to B'\), along with a continuous map
\(f_1\colon B\to B'\), a \emph{morphism of pre-mixed spaces over
\(f_1\)} is a continuous map \(f\colon V\to V'\) such that the diagram
\[
  \begin{CD}
    V @>f>> V'
  \\@V{\pi}VV @VV{\pi'}V
  \\B @>>{f_1}> B'
  \end{CD}
\] commutes and induces a \(\mathbb{C}\)-analytic map on each fibre.

A \emph{mixed space} is a pre-mixed space \(\pi\colon V\to B\) such that
every point \(y\in V\) admits a neighbourhood \(W\) in \(V\) that is
isomorphic as a pre-mixed space to an open subset of
\(B\times\mathbb{C}^n\), via an isomorphism over the identity. The
morphisms of mixed spaces are the same: mixed spaces form a \emph{full
subcategory}.

\subsection*{3. Deformations}\label{two-section-II.3}
\addcontentsline{toc}{subsection}{3. Deformations}

A mixed space \(\pi\colon V\to B\) is said to be \emph{proper} if \(B\)
is locally compact and the map \(\pi\) is proper (i.e.~the inverse image
of any compact subset is compact). If it is a mixed manifold, then we
can show that it is a fibred manifold that is locally trivial with
respect to the underlying \(\mathscr{C}^\infty\) structure, but the
previous talk shows that, in general, any two fibres are not isomorphic
as \(\mathbb{C}\)-analytic manifolds.

\begin{rmenv}{Definition}
Let \(V_0\) be a compact \(\mathbb{C}\)-analytic manifold, \(B\) a
locally compact space, and \(b_0\in B\). \oldpage{2-04} Then a
\emph{\(\mathbb{C}\)-analytic deformation of \(V_0\) over \((B,b_0)\)}
consists of a proper \(\mathbb{C}\)-analytic mixed space
\(\pi\colon V\to B\) along with an isomorphism of
\(\mathbb{C}\)-analytic manifolds \(i\colon V_0\to\pi^{-1}(b_0)\).

\end{rmenv}

The goal of this seminar is the study, at least local, and an attempt at
a classification of, \(\mathbb{C}\)-analytic deformations of a given
compact \(\mathbb{C}\)-analytic manifold \(V_0\).

\begin{rmenv}{Definition}
Let \(V_0\) be a compact \(\mathbb{C}\)-analytic manifold. A
\emph{\(\mathbb{C}\)-analytic deformation
\((\pi\colon V\to B,i\colon V_0\to V)\) of \(V_0\)} is said to be
\emph{locally complete} if, for any other deformation
\((\pi'\colon V'\to B',i'\colon V_0\to V')\) of \(V_0\), there exists a
neighbourhood \(B'_1\) of \(b'_0\) in \(B'\), an analytic map
\(f_1\colon B'_1\to B\) with \(f_1(b'_0)\to b_0\), and a morphism of
\(\mathbb{C}\)-analytic mixed spaces \(f\colon {\pi'}^{-1}(B'_1)\to V\)
over \(f_1\) such that \(f\circ i'=i\). The deformation is said to be
\emph{locally universal} is furthermore the germ of \(f_1\) at \(b'_0\)
is determined uniquely by this condition.

\end{rmenv}

It seems that every compact \(\mathbb{C}\)-analytic manifold \(V_0\)
admits a locally complete \(\mathbb{C}\)-analytic deformation, and a
locally universal one if the group of automorphisms of \(V_0\) is
discrete.

\section*{III. Vector fields}\label{two-section-III}
\addcontentsline{toc}{section}{III. Vector fields}

\subsection*{1. Study on models}\label{two-section-III.1}
\addcontentsline{toc}{subsection}{1. Study on models}

Let \(B\) be a space, \(U\) an object of \(\mathscr{S}_B\) (i.e.~an open
subset of \(B\times\mathbb{C}^n\)), \(b_0\) a point of \(B\), and set
\(U_0=\pi^{-1}(b_0)\).

A holomorphic field of tangent vectors on \(U_0\) (i.e.~a holomorphic
map from \(U_0\) to \(\mathbb{C}^n\)) is said to be a \emph{vertical
holomorphic field} on \(U_0\). A \emph{vertical holomorphic field on
\(U\)} is a continuous (resp. \ldots) map
\(\theta\colon U\to\mathbb{C}^n\) that induces a vertical holomorphic
field on each fibre \(U_x\). If \(f\colon U\to U'\) is an isomorphism in
\(\mathscr{S}_B\), then the \emph{transport \(f_*\theta\) of \(\theta\)
by \(f\)} is defined by \[
  f_*\theta(f(x,z)) = \mathrm{D}_2 f_{x,z}\cdot\theta(x,z)
\] where \(\mathrm{D}_2 f_{x,z}\) is the linear map from
\(\mathbb{C}^n\) to itself that is tangent to \(f_x\) at the point
\(z\in U_x\). This is again a vertical holomorphic field, since it
follows from a Cauchy integral that the matrix \(\mathrm{D}f_{x,z}\)
depends continuously on the pair \((x,z)\).

Now suppose that \(B\) is a \(\mathscr{C}^\infty\) manifold, just for
simplicity, and let \(T_0\) be the tangent space to \(B\) at \(b_0\).
\oldpage{2-05} A field of tangent vectors to \(U\) defined on \(U_0\),
i.e.~a map \(\omega\colon U_0\to T_0\times\mathbb{C}^n\), is said to be
a \emph{projectable holomorphic field} if
\(\omega(b_0,z)=(t_0,\theta(z))\) (where \(t_0\in T_0\) is a vector that
does not depend on \(z\), called the \emph{projection} of the field
\(\omega\)) and \(\theta(z)\) is a holomorphic vector field. If \(B\) is
a \(\mathbb{C}\)-analytic space, possibly with a singularity at \(b_0\),
then we give the same definition, but with \(T_0\) then being the
\emph{Zariski} tangent space to \(B\) at \(b_0\), i.e.~the dual of
\(\mathfrak{m}/\mathfrak{m}^2\), where \(\mathfrak{m}\) is the ideal of
germs at \(b_0\) of holomorphic functions on \(B\) that vanish at
\(b_0\).

If \(f\colon U\to U'\) is an isomorphism of
\(\mathscr{C}^\infty\mathscr{S}_B\) (resp. \ldots), then then transport
\(f_*\omega\) is defined by \[
  f_*\omega(f(b_0,z)) = \mathrm{D}f_{b_0,z}\omega(b_0,z)
\] where
\(\mathrm{D}f_{b_0,z}\colon T_0\times\mathbb{C}^n\to T_0\times\mathbb{C}^n\)
is now the linear map that is tangent to \(f\) at the point \((b_0,z)\).
This is a projectable holomorphic field. Indeed, the matrix
\(\mathrm{D}f_{b_0,z}\) can be written as \[
  \begin{pmatrix}
    I & 0
  \\\mathrm{D}_1f & \mathrm{D}_2f
  \end{pmatrix}
\] and \[
  \begin{aligned}
    \mathrm{D}_1f\colon T &\to \mathbb{C}^n
  \\\mathrm{D}_2f\colon \mathbb{C}^n &\to \mathbb{C}^n
  \end{aligned}
\] both depend holomorphically on \(z\) (for \(\mathrm{D}_1f\), this
follows from the fact that \(f_x\) is holomorphic for every \(x\)). By
setting \(f_*\omega(b_0,z')=(t_0,\theta'(z'))\), we have \[
  \begin{gathered}
    \theta'(z') = \mathrm{D}_1f_{b_0,z}(t_0) + \mathrm{D}_2f_{b_0,z}(\omega(z))
  \\\text{if }z'=f_{b_0}(z)
  \end{gathered}
\] which shows that \(f_*\omega\) is indeed a projectable holomorphic
field.

A \emph{projectable holomorphic field on \(U\)} is a
\(\mathscr{C}^\infty\) field of vectors tangent to \(U\) that induces a
projectable holomorphic field on each fibre.

\subsection*{2. Vector fields on a mixed
manifold}\label{two-section-III.2}
\addcontentsline{toc}{subsection}{2. Vector fields on a mixed manifold}

Let \(\pi\colon V\to B\) be a \((\mathscr{C}^\infty,\mathbb{C})\)-mixed
manifold (resp. \ldots, resp. a \(\mathbb{C}\)-analytic mixed space). By
transporting along the charts, we define the notions of

\begin{itemize}
\tightlist
\item
  vertical holomorphic fields on an open subset of a fibre;
\item
  vertical holomorphic fields on a open subset of \(V\);
\item
  projectable holomorphic fields on an open subset of a fibre; and
\item
  projectable holomorphic fields on an open subset of \(V\).
\end{itemize}

\oldpage{2-06}

Let \(\xi\) be a \(\mathscr{C}^\infty\) vector field (resp. \ldots) on
\(V\). By integrating \(\xi\), we obtain a \(\mathscr{C}^\infty\) map,
denoted by \(e^\xi\), from an open subset \(W\subset\mathbb{R}\times V\)
containing \(\{0\}\times V\) (resp. \(\mathbb{C}\)-analytic map from an
open subset \(W\subset\mathbb{C}\times V\)) to \(V\), characterised by

\begin{enumerate}
\def\labelenumi{\arabic{enumi}.}
\tightlist
\item
  \(e^\xi(t_1+t_2,y) = e^\xi(t_1,e^\xi(t_2,y))\), with the left-hand
  side being defined whenever the right-hand side is; and
\item
  \(\frac{\partial}{\partial t}e^\xi(t,y)|_{0,y} = \xi(y)\).
\end{enumerate}

Note that \(W\) is a mixed manifold over \(\mathbb{R}\times B\) (resp. a
mixed space over \(\mathbb{C}\times B\)).

\protect\phantomsection\label{two-proposition}
\begin{itenv}{Proposition}
For \(e^\xi\colon W\to V\) to be a morphism of mixed spaces over the
projection \(\mathbb{R}\times B\to B\), it is necessary and sufficient
for \(\xi\) to be a vertical holomorphic field. For
\(e^\xi\colon W\to V\) to be a morphism of mixed spaces over a map from
an open subset of \(\mathbb{R}\times B\) containing \(\{0\}\times B\) to
\(B\), it is necessary and sufficient for \(\xi\) to be a projectable
holomorphic field.

\end{itenv}

The proof is left to the reader.

\section*{IV. The Spencer--Kodaira map}\label{two-section-IV}
\addcontentsline{toc}{section}{IV. The Spencer--Kodaira map}

Let \(\pi\colon V\to B\) be a mixed manifold (resp. a
\(\mathbb{C}\)-analytic mixed space), \(b\in B\), and
\(V_0=\pi^{-1}(b_0)\). Let \(T_0\) be the tangent space to \(B\) at
\(b_0\) (resp. the Zariski tangent space). We introduce the following
sheaves on \(V_0\):

\begin{itemize}
\tightlist
\item
  \(\Theta_0\): the sheaf of germs of vertical holomorphic fields on
  \(V_0\) ;
\item
  \(\Pi_0\): the sheaf of germs of locally projectable holomorphic
  fields on \(V_0\) ; and
\item
  \(\Lambda_0\): the sheaf \(\pi^*T_0\), i.e.~the sheaf of germs of
  locally constant maps from \(V_0\) to \(T_0\).
\end{itemize}

We have an exact sequence of sheaves on \(V_0\) \[
  0 \to \Theta_0 \to \Pi_0 \to \Lambda_0 \to 0
\] that gives rise to the long exact sequence in cohomology \[
  \ldots \to \mathrm{H}^0(V_0;\Pi_0) \to \mathrm{H}^0(V_0;\Lambda_0) \xrightarrow{\delta} \mathrm{H}^1(V_0;\Theta_0) \to \ldots.
\] We also have a canonical map \oldpage{2-07} \[
  \iota\colon T_0 \to \mathrm{H}^0(V_0;\Lambda_0)
\] that is injective if \(V_0\) is non-empty, and surjective if \(V_0\)
is connected.

\protect\phantomsection\label{two-definition}
\begin{rmenv}{Definition}
The \emph{Spencer--Kodaira map} is the composition \[
  \rho_0 = \delta\circ\iota\colon T_0 \to \mathrm{H}^1(V_0;\Theta_0).
\]

\end{rmenv}

This map is an essential tool in the local study of deformations of
\(\mathbb{C}\)-analytic varieties. Note that \(\Theta_0\) is exactly the
sheaf of germs of holomorphic fields of tangent vectors to \(V_0\), and
thus depends only on \(V_0\), while \(T_0\) depends only on the base.
Also, \(\Theta_0\) is a coherent analytic sheaf on \(V_0\), and, if
\(V_0\) is compact, then \(\mathrm{H}^1(V_0;\Theta_0)\) is a
finite-dimensional vector space over \(\mathbb{C}\) \autocite{2-1}. We
thus see that, in this case (which is the only case where we can say
anything non-trivial), \(\rho_0\) might be possible to calculate.

It is clear that, if the given mixed manifold is trivial (i.e.~if
\(V=B\times V_0\), with \(\pi\) being the projection to \(B\)), then the
map \(\rho_0\) is zero. The next talk aims to show that, in a certain
sense, \(\rho\) indicates the non-triviality of \(V\) in a neighbourhood
of \(V_0\).


% Bibliography
\nocite{*}
\printbibliography[heading=bibintoc,title=Bibliography]

\end{document}
